% generated by roam.py handbook export -- не редактировать, править заметки
% 20260922181524-оценка_нормы_гессиана_сверху_для_функции_с_липшицевым_градиентом.org (body 9d5497de53e8)

\begin{theorem}[{Оценка нормы гессиана сверху для функции с липшицевым градиентом}]\label[theorem]{note:DAA0DAE4-4B48-4A2F-9459-54790B6B129F}
  Пусть $E \subset \R^n$~--- открытое множество, $f: E \rightarrow \R$, $f \in C^{2,1}_L$: $f$ дважды непрерывно дифференцируема и ее \hyperref[note:7DD1C7AB-1A49-4675-9F21-5A16C717F670]{градиент липшицев} с константой $L$.

  Тогда  $\| \nabla^2 f(x) \| \le L\ \forall x \in E$, где $\| \cdot \|$~--- спектральная норма
\end{theorem}

\begin{remark}
  Обратное~--- \hyperref[note:4B6D53BA-E599-4DEA-A557-550BD32042BA]{связь спектральной нормы гессиана и липшицевости градиента}.
\end{remark}

\begin{proof}
  Возьмем $h \in \R^n$ и $t > 0$. $\nabla f$~--- дифференцируема $\Rightarrow \nabla f(x + th) - \nabla f(x) = \nabla^2 f(x)th + o(th)$ \\
  Тогда нормы двух частей тоже равны:
  \[
  \|  \nabla^2 f(x)th + o(th) \| = \| \nabla f(x + th) - \nabla f(x) \| \le L \| x + th - x \|
  \]
  По липшицевости градиента.
  Давайте разделим обе части неравенства на $t$, и устремим $t \rightarrow 0+$
  \[
  \|  \nabla^2 f(x)h + \frac{o(th)}{t}  \| \le L \| h \|
  \]
  Мы знаем, что $\frac{o(th)}{t} \rightarrow 0$. Тогда по предельному переходу
  \[
  \|  \nabla^2 f(x)h  \| \le L \| h \|
  \]
  Так как $h$ произвольный, по определению спектральной нормы $\| \nabla^2 f(x) \| = \sup\limits_{h \ne 0} \frac{\| \nabla^2 f(x) h \|}{\| h \|} \le L$.
\end{proof}
